\documentclass[12pt]{article}
\usepackage[utf8]{inputenc}
\usepackage[T1]{fontenc}
\usepackage{hyperref}
\title{Yadgha Language Resources for K-12 Curriculum Development}
\author{Compiled by Claude Sonnet 5.5 (claude-sonnet-5-5) for Andr\'as Kornai}
\date{}
\begin{document}
\maketitle

\section{Introduction}
Yadgha (Yidgha, ISO 639-3 ydg) is an Eastern Iranian language of the upper Lotkoh (Lutkuh) Valley in Chitral, Khyber Pakhtunkhwa, Pakistan. It is closest to Munji, which is spoken across the border in Afghanistan. Estimates give roughly 6,000--7,000 speakers. Speakers are increasingly shifting to Khowar, and UNESCO classes the language as endangered.

\section{Primary Educational Resources}
No school textbooks were found. The only publications reported are those of the Forum for Language Initiatives (FLI) 2016 project, which recorded ten folktales with Urdu translations and produced the first publications in Yadgha.
\begin{itemize}
\item FLI, ``Yidgha will live a long life''. \url{https://fli-online.org/yidgha-will-live-a-long-life/}
\end{itemize}

\section{Linguistic and Reference Materials}
\begin{itemize}
\item FLI, ``Yidgha: an endangered language in Chitral'' (report). \url{https://www.fli-online.org/documents/languages/yidgha/yidgha-report.htm}
\item Wikipedia, ``Yidgha language'', which mentions Georg Morgenstierne's 1938 description of Yidgha and Munji with a glossary. This is the main lexical resource, and no modern dictionary was found. \url{https://en.wikipedia.org/wiki/Yidgha_language}
\item Omniglot, Yidgha, on the writing system. \url{https://www.omniglot.com/writing/yidgha.htm}
\end{itemize}

\section{Religious and Community Resources}
\begin{itemize}
\item FLI, Yidgha page. \url{https://fli-online.org/yidgha/}
\end{itemize}

\section{Digital and Audio Resources}
\begin{itemize}
\item Endangered Languages Project, Yidgha entry. \url{https://endangeredlanguages.com/lang/3210}
\end{itemize}
No audio files were directly located in the search results.

\section{Curriculum Development Context}
There is no official orthography. Some people use the Arabic-based Urdu script with extra letters, and FLI later helped the community adopt a writing system. A curriculum would build on the FLI folktales and orthography work, with community leadership. It would also address the shift to Khowar and Urdu.
\end{document}